\documentclass[11pt]{article}
\usepackage[utf8]{inputenc}
\usepackage[margin=1in]{geometry}
\usepackage{natbib}
\usepackage{booktabs}
\usepackage{amsmath}
\usepackage{graphicx}
\usepackage{hyperref}
\usepackage{xcolor}

\title{Design Principles for Inflammasome Inhibition by Pyrin-Only Proteins}

\author{
  Author One$^{1}$, Author Two$^{1}$, Author Three$^{2}$, Author Four$^{1,*}$ \\
  $^{1}$Department of Immunology, University A \\
  $^{2}$Department of Biochemistry, University B \\
  $^{*}$Corresponding author
}

\date{}

\begin{document}
\maketitle

\begin{abstract}
Inflammasomes are supramolecular innate immune signaling platforms whose dysregulation underlies autoinflammatory diseases and cancer. Pyrin-only proteins (POPs) are endogenous single-domain inhibitors of inflammasome assembly, yet their precise target specificities and inhibitory mechanisms remain poorly characterized. Here we use a combination of Rosetta computational modeling, in vitro filament assembly assays, and cell-based inflammasome activation assays to systematically map the target specificity and mechanisms of POP1 and POP2. Contrary to a previous report, we find that POP1 is a poor inhibitor of the central adaptor ASC. Instead, POP1 blocks the assembly of upstream receptor pyrin domain (PYD) filaments from AIM2, IFI16, NLRP3, and NLRP6. POP2, by contrast, directly suppresses ASC nucleation and additionally inhibits filament elongation of AIM2 and NLRP6. Rosetta modeling predicts differential binding modes that explain these distinct specificities, and cell-based assays confirm the full inhibition profile. Our results reveal a division of labor in which POP1 intercepts upstream receptor filament formation while POP2 targets the adaptor layer, providing complementary layers of inflammasome regulation. These findings establish design principles for how small PYD-only proteins achieve selective inflammasome inhibition.
\end{abstract}

\section{Introduction}

Inflammasomes are multiprotein complexes that assemble in the cytosol in response to pathogen- and damage-associated molecular patterns (PAMPs and DAMPs) \citep{martinon2002,broz2016}. Canonical inflammasome activation involves the oligomerization of sensor proteins (e.g., NLRP3, AIM2, NLRP6, IFI16) through their pyrin domains (PYDs), recruitment of the adaptor protein ASC via homotypic PYD--PYD interactions, and subsequent caspase-1 activation leading to cytokine maturation and pyroptotic cell death \citep{schroder2010,lu2014}.

The assembly of inflammasome signaling platforms proceeds through a nucleated polymerization mechanism in which sensor PYD filaments nucleate ASC PYD filament assembly, which in turn nucleates caspase-1 filaments in a prion-like amplification cascade \citep{lu2014,cai2014}. This architecture provides signal amplification but also poses a risk of aberrant activation. Indeed, dysregulated inflammasome signaling contributes to numerous autoinflammatory diseases, gout, atherosclerosis, neurodegeneration, and cancer \citep{masters2009,duewell2010,heneka2013}.

Endogenous mechanisms that restrain inflammasome activation are therefore critical for immune homeostasis. Pyrin-only proteins (POPs) were identified as a class of endogenous inflammasome inhibitors that consist solely of a PYD and lack effector domains \citep{stehlik2007,dorfleutner2007}. Two human POPs---POP1 and POP2---have been characterized, and they are thought to inhibit inflammasome assembly by competing with productive PYD--PYD interactions \citep{khare2014,ratsimandresy2017}. However, the precise target specificity of each POP has been unclear. A previous study reported that POP1 directly inhibits ASC \citep{stehlik2003}, but this finding has not been thoroughly validated with modern structural and biochemical approaches.

Here we undertake a systematic analysis of POP1 and POP2 target specificity using computational modeling with Rosetta, quantitative in vitro filament assembly assays, and cell-based inflammasome activation experiments. Our results revise the current model of POP function and establish design principles for selective inflammasome inhibition by PYD-only proteins.

\section{Related Work}

\subsection{Inflammasome Assembly Mechanisms}

The structural basis of inflammasome assembly has been elucidated through cryo-EM studies of PYD filaments. The ASC PYD filament structure revealed a helical assembly mediated by three types of asymmetric PYD--PYD interactions \citep{lu2014}. AIM2 PYD filaments were shown to nucleate ASC assembly through direct PYD--PYD contacts \citep{lu2015}. NLRP3 and NLRP6 PYD filaments have been characterized more recently, revealing conserved but non-identical interaction surfaces \citep{sborgi2015,sharif2019}.

\subsection{Endogenous Inflammasome Regulators}

Multiple endogenous regulators of inflammasome signaling have been identified, including POPs, CARD-only proteins (COPs), and regulatory splice variants \citep{stehlik2007,dorfleutner2007}. POPs are small ($\sim$10~kDa) proteins that contain only a PYD and are proposed to act as decoys or competitive inhibitors \citep{khare2014}. POP3 (a third family member) was shown to inhibit AIM2 and IFI16 inflammasomes \citep{khare2014}. However, the mechanistic distinction between POP1 and POP2 has remained a significant knowledge gap.

\subsection{Computational Modeling of PYD Interactions}

Rosetta-based computational modeling has been applied to predict protein--protein interactions in the inflammasome field, including PYD--PYD binding interfaces \citep{de2012,sborgi2015}. The RosettaDock protocol enables the prediction of binding poses and interface energetics for homotypic PYD interactions \citep{gray2003,chaudhury2011}. These computational predictions can be validated by mutagenesis and functional assays, providing a powerful approach for mapping interaction specificity.

\section{Methods}

\subsection{Rosetta Computational Modeling}

POP1 and POP2 structures were modeled using RosettaCM comparative modeling \citep{song2013}. Docking of POPs to PYD filament surfaces (ASC, AIM2, IFI16, NLRP3, NLRP6) was performed using RosettaDock 4.0 \citep{gray2003,chaudhury2011}. Interface energies were scored using the Rosetta energy function (REF2015) \citep{alford2017}. For each POP--PYD pair, 10,000 docking trajectories were generated and the top 100 by interface score were analyzed for binding mode classification.

\subsection{In Vitro Filament Assembly Assays}

Recombinant PYD proteins (ASC PYD, AIM2 PYD, IFI16 PYD, NLRP3 PYD, NLRP6 PYD) were expressed in \textit{E.~coli} and purified under denaturing conditions. Filament assembly was initiated by rapid dilution into physiological buffer and monitored by dynamic light scattering and negative-stain electron microscopy \citep{lu2014}. POP1 or POP2 was added at varying molar ratios (0.5:1 to 5:1 POP:PYD) to assess inhibition of filament formation.

\subsection{Cell-Based Inflammasome Activation Assays}

HEK293T cells were reconstituted with individual inflammasome components (sensor + ASC + caspase-1) and stimulated with appropriate ligands. POP1 or POP2 was co-expressed, and inflammasome activation was measured by caspase-1 cleavage (western blot), ASC speck formation (fluorescence microscopy), and IL-1$\beta$ secretion (ELISA) \citep{stutz2013}.

\section{Results}

\subsection{POP1 Is a Poor Inhibitor of ASC}

Contrary to a previous report \citep{stehlik2003}, we find that POP1 does not effectively inhibit ASC PYD filament assembly or ASC-mediated inflammasome signaling. In vitro, POP1 at a 5:1 molar ratio produced minimal inhibition of ASC PYD polymerization. In cell-based assays, POP1 co-expression did not significantly reduce ASC speck formation upon NLRP3 stimulation when POP1 was the only inhibitor present.

\subsection{POP1 Inhibits Upstream Receptor PYD Filaments}

Instead of targeting ASC, POP1 potently inhibited the assembly of upstream receptor PYD filaments (Table~\ref{tab:pop1}). POP1 blocked filament formation of AIM2 PYD, IFI16 PYD, NLRP3 PYD, and NLRP6 PYD in vitro. Cell-based assays confirmed that POP1 reduced caspase-1 activation downstream of all four receptors, consistent with inhibition at the receptor filament level.

\begin{table}[ht]
\centering
\caption{POP1 target specificity: revised findings.}
\label{tab:pop1}
\begin{tabular}{llll}
\toprule
\textbf{Target} & \textbf{POP1 Inhibition} & \textbf{Previous Report} & \textbf{This Study} \\
\midrule
ASC (adaptor) & Poor & Strong inhibitor & Contradicted \\
AIM2 PYD filament & Assembly blocked & Not tested & New finding \\
IFI16 PYD filament & Assembly blocked & Not tested & New finding \\
NLRP3 PYD filament & Assembly blocked & Not tested & New finding \\
NLRP6 PYD filament & Assembly blocked & Not tested & New finding \\
\bottomrule
\end{tabular}
\end{table}

\subsection{POP2 Targets ASC Nucleation and Receptor Filament Elongation}

POP2 exhibited a distinct and dual mechanism of inhibition (Table~\ref{tab:pop2}). First, POP2 directly suppressed ASC PYD nucleation, blocking the initial polymerization event. Second, POP2 inhibited filament elongation of AIM2 and NLRP6 PYDs, acting as a capping factor that terminates growing filaments. POP2 did not significantly affect NLRP3 or IFI16 filament elongation, indicating receptor-specific capping activity.

\begin{table}[ht]
\centering
\caption{POP2 target specificity and mechanisms.}
\label{tab:pop2}
\begin{tabular}{lll}
\toprule
\textbf{Target} & \textbf{POP2 Activity} & \textbf{Mechanism} \\
\midrule
ASC nucleation & Direct suppression & Blocks initial nucleation \\
AIM2 filament elongation & Inhibits & Caps growing filament \\
NLRP6 filament elongation & Inhibits & Caps growing filament \\
\bottomrule
\end{tabular}
\end{table}

\subsection{Comparative Analysis of POP1 and POP2}

The systematic comparison reveals a clear division of labor between POP1 and POP2 (Table~\ref{tab:comparison}). POP1 acts primarily at the receptor level, blocking the initial assembly of sensor PYD filaments across a broad range of receptors (AIM2, IFI16, NLRP3, NLRP6). POP2 acts primarily at the adaptor level, directly inhibiting ASC nucleation, while also providing filament-capping activity at a subset of receptors (AIM2, NLRP6). This complementary targeting ensures that inflammasome assembly can be inhibited at multiple steps in the signaling cascade.

\begin{table}[ht]
\centering
\caption{Comparison of POP1 and POP2 mechanisms and specificities.}
\label{tab:comparison}
\begin{tabular}{lll}
\toprule
\textbf{Feature} & \textbf{POP1} & \textbf{POP2} \\
\midrule
Primary target & Upstream receptor PYD filaments & ASC nucleation + receptor filaments \\
ASC inhibition & Poor & Strong (direct) \\
Receptor inhibition & Assembly blockade & Elongation blockade \\
Receptors affected & AIM2, IFI16, NLRP3, NLRP6 & AIM2, NLRP6 \\
\bottomrule
\end{tabular}
\end{table}

\subsection{Rosetta Modeling Predicts Differential Binding Modes}

Rosetta docking simulations predicted that POP1 preferentially binds receptor PYD surfaces with favorable interface energies, while showing poor complementarity with the ASC PYD surface. Conversely, POP2 docking predicted favorable binding to both ASC PYD (at the nucleation interface) and to the filament-capping sites of AIM2 and NLRP6 PYDs. These computational predictions are consistent with the experimentally observed specificity profiles and suggest that electrostatic complementarity at the binding interface is a major determinant of POP target selectivity.

\subsection{Validation by Cell-Based Assays}

Cell-based inflammasome activation assays confirmed the specificity map derived from in vitro and computational approaches (Table~\ref{tab:methods}). POP1 expression reduced caspase-1 activation downstream of AIM2, IFI16, NLRP3, and NLRP6 but did not significantly affect ASC-dependent signaling when receptor filament formation was bypassed. POP2 expression suppressed ASC-mediated caspase-1 activation and additionally reduced signaling through AIM2 and NLRP6 pathways.

\begin{table}[ht]
\centering
\caption{Summary of experimental approaches and key results.}
\label{tab:methods}
\begin{tabular}{lll}
\toprule
\textbf{Approach} & \textbf{Application} & \textbf{Key Result} \\
\midrule
Rosetta in silico & POP--PYD binding prediction & Differential binding modes \\
In vitro filament assay & Filament inhibition & POP1/POP2 inhibitory profiles \\
In cellulo assay & Functional validation & Confirmed selective inhibition \\
\bottomrule
\end{tabular}
\end{table}

\section{Discussion}

Our systematic analysis of POP target specificity revises the current model of inflammasome regulation by pyrin-only proteins and establishes several design principles for how these small single-domain proteins achieve selective inhibition of a complex signaling cascade.

The most significant finding is the correction of POP1's previously reported target. The prior identification of POP1 as an ASC inhibitor \citep{stehlik2003} was based on overexpression co-immunoprecipitation experiments that may have been confounded by the high local concentrations of PYD domains. Our quantitative in vitro and cell-based assays, combined with Rosetta modeling, consistently demonstrate that POP1 preferentially targets upstream receptor PYD filaments rather than ASC. This revision has important implications for understanding how endogenous inflammasome regulation is organized.

The complementary targeting by POP1 and POP2 creates a layered defense against aberrant inflammasome activation. POP1 acts as a broad-spectrum receptor filament inhibitor, targeting the initial assembly step for four different inflammasome sensors. POP2 provides a second checkpoint at the adaptor level by directly inhibiting ASC nucleation, while also serving as a filament-capping factor for a subset of receptors. This architecture ensures that even if receptor filament formation proceeds past the POP1 checkpoint, POP2 can still prevent signal amplification through ASC.

The receptor selectivity of POP2's capping activity---affecting AIM2 and NLRP6 but not NLRP3 or IFI16---likely reflects structural differences in the filament-capping interfaces of these receptors. The Rosetta modeling suggests that the electrostatic surface of POP2 is particularly complementary to the capping sites of AIM2 and NLRP6, consistent with the known sequence divergence among inflammasome sensor PYDs \citep{lu2015,sborgi2015}.

These findings have implications for therapeutic targeting of inflammasomes. Current drug development efforts focus primarily on inhibiting NLRP3 directly \citep{coll2015,mangan2018}. The design principles revealed by POP specificity suggest that targeting receptor PYD filament assembly (the POP1 mechanism) or ASC nucleation (the POP2 mechanism) could provide alternative therapeutic strategies with distinct selectivity profiles.

\section{Conclusion}

We have systematically mapped the target specificities and inhibitory mechanisms of POP1 and POP2 using an integrated computational, biochemical, and cellular approach. POP1 is a broad-spectrum inhibitor of upstream receptor PYD filaments (AIM2, IFI16, NLRP3, NLRP6) but a poor ASC inhibitor, contradicting previous reports. POP2 directly suppresses ASC nucleation and caps AIM2 and NLRP6 filament elongation. This division of labor provides complementary layers of inflammasome regulation and establishes design principles for selective inflammasome inhibition by PYD-only proteins.

\bibliographystyle{plainnat}
\bibliography{references}

\end{document}
